El muestreo permite estudiar poblaciones mediante una parte de sus unidades, pero exige mantener una correspondencia explícita entre la población delimitada y la evidencia realmente obtenida. La diferencia entre parámetro y estadístico aclara qué se desea conocer y qué se observa; la diferencia entre variación de muestreo y sesgo aclara qué puede mejorar al aumentar el número de observaciones y qué requiere cambiar el procedimiento de recolección \citep{wackerly2010,bradley2021}.

La selección aleatoria simple, la estratificación, el muestreo sistemático y los conglomerados responden a necesidades distintas de cobertura, organización y costo. El criterio administrativo debe ser elegir el diseño que corresponda a la población y a la decisión, en lugar de aplicar una técnica solo porque resulta conocida o accesible. La presencia de grupos en una muestra merece evaluarse junto con la manera en que sus integrantes fueron elegidos y con la respuesta que efectivamente se obtuvo \citep{holmesBusiness2023}.

Considero que una práctica responsable debe priorizar la calidad del marco y la transparencia de la selección antes de interpretar la cantidad de registros como evidencia suficiente. En una encuesta al personal, ello implica incluir turnos y áreas con distintas condiciones de acceso; en una revisión de transacciones, delimitar el periodo, los registros elegibles y las reglas de selección. Si la cobertura excluye sistemáticamente un segmento, aumentar la muestra puede dar una impresión de precisión sin resolver el problema de fondo \citep{bradley2021}.

El tamaño se planifica con un nivel de confianza, una tolerancia y supuestos coherentes con la variable. Sin embargo, esa planificación debe acompañarse de seguimiento de la no respuesta, revisión de medición y comunicación de límites. Los resultados pueden orientar recursos o decisiones laborales cuando se respeta su alcance; extrapolarlos a otras poblaciones o atribuirles efectos causales requiere evidencia adicional \citep{bloom2024}.

Una conclusión útil para la administración vincula selección, medida y decisión: indica a quiénes describe, cuánta incertidumbre admite y qué limitaciones deben considerarse antes de actuar. Esa disciplina permite utilizar el muestreo como fundamento de decisiones proporcionales a la evidencia disponible.
